\section[Further exact source information on the three broad C1 rows]{Further exact source information on the three broad C1 rows}
\label{sec:c1-residual-source-implications}

This section records source implications of (NC1)--(NC7), not further
all-tests certificates. All quantities are in the coefficient field.
Write $f\asymp g$ for two strictly positive quantities whose quotient
and reciprocal quotient are bounded. Set
\[
             L=x\zeta,\qquad S=1+L+y,\qquad \rho=r+\eta.
\]
The native signs and (NC5) imply the uniform estimates
\[
 S\le P_d\le6S,\qquad
 \frac{1+L}{54}\le P_4\le3(1+L),\qquad
 P_\alpha\asymp h(1+y)+\zeta M.                 \tag{NI1}
\]
Indeed $rx^2\zeta<x$ follows from $Q>0$, and
$\eta x^2h\theta=x\theta(1-E)\le1$, giving the $P_d$ bounds
from (NC2). For $P_4$, if $L\le2$ use $E\ge1/18$; if $L\ge2$,
then $r<1/L\le1/2$ and $P_4\ge L/2$. Its upper bound follows
immediately from (NC2). The last comparison uses $1\le R\le2$
and $1/18\le E<1$.
Consequently the unchanged original source satisfies
\begin{align*}
 C/d&\asymp\frac{\rho}{(1-r)S},&
 A/d&\asymp\frac{\eta(1+L)}{rS},\\
 m/d&\asymp\frac{x\rho h(1+y)}{M\zeta S},&
 \alpha/d&\asymp
 \frac{x^2\zeta\rho\{h(1+y)+\zeta M\}}{QS}.
 \tag{NI2}
\end{align*}
No opening has been bounded away from zero in these statements.
In particular, $E$ and $P_4$ do not supply additional hidden small
source denominators on C1.1, C1.4 or C1.5.

\paragraph{C1.1: two unavoidable large weights and its actual angle.}
Here $x,\zeta$ are infinitesimal and $\eta$ is unbounded.
The inherited original-$D$ cap gives $q=x\eta$ bounded with
residue at most $1/2$. Hence $y=q\zeta$ is infinitesimal:
the actual angle type is necessarily $(I,I)$, independently of $e$.
The exact formulas give
\[
 C/d\asymp\frac{\eta}{1-r},\qquad
 A/d\asymp\frac{\eta}{r},\qquad
 m/d\asymp\frac{e^2}{M},\qquad
 \alpha/d\asymp xq\zeta,\qquad
 \beta/d+\frac12\asymp\frac{e^2}{t}.             \tag{NI3}
\]
Thus $C,A$ are unbounded relative to $d$, and $\alpha$ is weak.
The inherited bounded polynomial formula for $B_5$ gives $D/d=O(1/\eta)$,
so $D$ is weak too, including its literal zero. In (NI3), $\beta\ge0$
forces $t=O(e^2)$. Therefore when $e$ is infinitesimal, $t$ is
infinitesimal, $M$ is a unit and $m$ is weak. When $e$ is a unit,
$m$ is strong; when $e$ is unbounded, both $m,\beta$ are unbounded.
These are restrictions on the original C1.1 row, not new admissions.
Its positive infinitesimal $\delta_{\rm tot}$ condition is automatic
because $C/d,A/d$ are unbounded.

\paragraph{C1.4: the natural scale is $L=x\zeta$.}
Now $x$ is infinitesimal and $\zeta$ is noninfinitesimal.
The following implications follow uniformly through all retained
coefficient boundaries.
\begin{enumerate}
\item If $L$ is unbounded, $Q>0$ forces $r$ infinitesimal, and
$1-r=a\eta<2\eta$ makes $\eta$ noninfinitesimal. Hence $y=L\eta$
is unbounded: the actual type is $(I,H)$. Here $\zeta$ is unbounded,
$t=\theta/\zeta$ is infinitesimal, and $M$ is a unit. Since
$h y/\zeta=x\eta h<1$, (NC2) gives $P_\alpha\asymp\zeta$.
Equations (NI1)--(NI2) then yield
\[
 C/d\asymp L^{-1},\qquad
 A/d\asymp r^{-1},\qquad
 \alpha/d\asymp L/Q,\qquad
 m/d\asymp (x\eta h)/\zeta\le O(\zeta^{-1}).
 \tag{NI4}
\]
In particular $C,m$ are weak, $A,\alpha$ unbounded, and
$(m/d)/(C/d)=O(x)$ is infinitesimal. The original total-scale
intersection is automatic.

\item If $L$ is a bounded unit, $\zeta$ is unbounded, $M$ is a
unit, $P_\alpha\asymp\zeta$, and $S\asymp1+\eta$. The identity
$1=r+a\eta$ gives $\rho/(1+\eta)$ a positive bounded unit.
Thus
\[
 C/d\asymp(1-r)^{-1},\qquad
 \alpha/d\asymp Q^{-1},\qquad
 A/d\asymp\frac{\eta}{r(1+\eta)},\qquad
 m/d=O(x).
 \tag{NI5}
\]
So $C,\alpha$ are strong and $m$ weak. Moreover $A$ is weak
exactly when $\eta$ is infinitesimal; that case has $r$ a unit and
$D/d$ unbounded, by $B_5\sim1$ in (NC3). The actual Bob angle is
infinitesimal, a unit, or unbounded exactly as $\eta$ is, because
$y=L\eta$. No restriction on $D,\beta$ is added in the other cases.

\item If $L$ is infinitesimal and $\eta$ is bounded, then $y$ is
infinitesimal. The $B_\beta\ge0$ equation bounds
$\theta\le Kx\rho G/P_d=O(x+L)$, so $h,M$ are units.
Now (NI2) gives $C$ strong, $m,\alpha$ weak, and $A$ weak exactly
when $\eta$ is infinitesimal. In that last case $D/d$ is unbounded,
as in the preceding item. The actual type is $(I,I)$.

\item If $L$ is infinitesimal and $\eta$ is unbounded, then
\[
 C/d\asymp\frac{\eta}{(1-r)(1+y)},\qquad
 A/d\asymp\frac{\eta}{r(1+y)}.                    \tag{NI6}
\]
Both are unbounded since $\eta/(1+y)=(\eta^{-1}+L)^{-1}$.
Also $Q$ is a unit and (NI2) bounds
\[
 \alpha/d\le O\!\left(x^2\zeta\eta h+
              \frac{x^2\zeta^2\eta}{1+y}\right)=O(L),
\]
so $\alpha$ is weak. Finally $J\sim1$ and
$P_d/(1+y)\sim1$ directly in (NC2); (NC4) gives
$D/d+1/2\sim1/2$, hence $D$ weak. This extends the
$C,A$-large/$D,\alpha$-weak source pattern of C1.1 into C1.4,
without claiming that its existing conductors apply. Here $y$ may
be infinitesimal, a unit, or unbounded. The total-scale intersection
is again automatic.
\end{enumerate}
These four alternatives exhaust C1.4, but they are necessary source
reductions only. In particular no entire C1.4 leaf is closed.

\paragraph{C1.5 after the actual unit/unit admission.}
Here $k$ is a unit. The new theorem removes actual $(U,U)$; the
inherited both-$\eta,\zeta$-infinitesimal theorem is also subtracted.
Two useful remaining source types can be specified more tightly.

If $v_c$ is infinitesimal, $\eta\zeta$ is infinitesimal. Neither
$\eta$ nor $\zeta$ can be unbounded: $\zeta$ infinitesimal forces
$\theta$ infinitesimal and excludes unbounded $\eta$ by $E>0$;
$\eta$ infinitesimal gives $r\sim1$ and excludes unbounded $\zeta$
by $Q>0$. Thus exactly one of $\eta,\zeta$ is a unit in this residual
subcase. The exact stabilizer (AS5) exchanges these two alternatives.
It keeps $k$ a bounded unit, although a second orientation need not
keep $k\le1$. Choose the representative $\eta$ a unit and
$\zeta$ infinitesimal. Then
\[
 m/d,\beta/d\text{ are unbounded},\qquad
 C/d,A/d\text{ are strong},\qquad \alpha/d\text{ is weak}.
 \tag{NI7}
\]
Indeed $\theta=t\zeta$ is infinitesimal, $E$ is a unit,
$P_d,P_4,P_\alpha,r+\eta$ are units, and (NC4) gives
$m/d\asymp1/(M\zeta)$ and $\beta/d+1/2\asymp1/\theta$.
The reflected representative also has a unit $E$: its exact value
is $E'=h(1+y)/G$ in the original coordinates, with $h$ a unit by
the same $\theta=1$ exclusion used in the ordinary-angle proof.
There is no conclusion about the size of $D/d$ in (NI7).

If $v_c$ is unbounded, both $\eta,\zeta$ are noninfinitesimal and
at least one is unbounded. To see the first assertion, an infinitesimal
$\zeta$ would force $h\sim1$ and $\eta$ bounded by $E>0$;
an infinitesimal $\eta$ would force $r\sim1$ and $\zeta$ bounded
by $Q>0$. Two rate regimes are then useful.
\begin{itemize}
\item If $\eta$ is unbounded and $\zeta$ a unit, then $h$ is
infinitesimal. The exact alternative numerator identity
\[
 B_\beta=(x\eta h)\{K(1+\zeta)-x\theta^2\}
             +xrGT-\theta(1+2xh+x\zeta)             \tag{NI8}
\]
shows that $x\eta h$ is a unit: otherwise its residue at
$\theta=1$ would be $-1-x\zeta(1-r)<0$.
Thus $P_\alpha$ is a unit and $P_d\asymp\eta$.
Equations (NC4) give
\[
 C/d\asymp(1-r)^{-1},\quad A/d\asymp r^{-1},\quad
 m/d\asymp M^{-1},\quad\alpha/d\asymp Q^{-1},
 \quad D/d,\beta/d\text{ weak}.                    \tag{NI9}
\]
For the last assertion use $J\sim1$, $\theta\sim1$ and
$P_d/(Ky)\sim1$ in the original $D/d+1/2,\beta/d+1/2$
formulas. The total-scale condition $\delta_{\rm tot}$ infinitesimal
is therefore equivalent on this typed branch to at least one of
$r,1-r,M,Q$ being infinitesimal. If all four are units, the inherited
positive-total-scale theorem already applies.

\item If $\zeta$ is unbounded, $r$ is infinitesimal and $M$ a
unit. The noninfinitesimal $\eta$ obeys $\eta h\asymp1$.
For if $\theta$ is infinitesimal this follows from $E>0$ and
$h\sim1$; otherwise divide (NI8) by $\zeta$, use
$xrG/\zeta=O(r)$, and obtain that $x\eta h$ has a positive-unit
lower bound. Therefore (NC2) and (NC4) give
\[
 C/d\asymp\zeta^{-1},\quad m/d\asymp\zeta^{-1},\quad
 A/d\asymp r^{-1},\quad\alpha/d\asymp\zeta/Q.
 \tag{NI10}
\]
Thus $C,m$ are weak and $A,\alpha$ unbounded. No size restriction
on $D,\beta$ is inferred. This branch already has infinitesimal
$\delta_{\rm tot}$.
\end{itemize}
All implications retain the original physicality and exact source
openings. They identify the missing source scales for further full-parent
proofs; only the actual $(U,U)$ subcase has received a new all-tests
admission in this section and its companion.
